% Draft for IEEE CSF, translated from the repository documentation (README.md,
% docs/*.md).  Class: IEEEtran in conference mode, as CSF requires.
% TODO: check the current CSF call for the page limit and whether submissions
% must be anonymised (if so, replace the author block below).
\PassOptionsToPackage{obeyspaces,spaces}{url}
\documentclass[conference]{IEEEtran}

\usepackage[T1]{fontenc}
\usepackage[utf8]{inputenc}
\usepackage{amsmath,amssymb}
\usepackage{amsthm}
\usepackage{booktabs}
\usepackage{tabularx}
\newcolumntype{L}{>{\raggedright\arraybackslash}X}
\usepackage{xcolor}
\usepackage{listings}
\usepackage{hyperref}
\hypersetup{colorlinks,linkcolor=black,citecolor=black,urlcolor=blue!50!black}
\usepackage[capitalise,nameinlink]{cleveref}

% ---------------------------------------------------------------- listings
\lstdefinelanguage{Coq}{
  morekeywords={match,with,end,if,then,else,let,in,fun,forall,exists,
    Definition,Lemma,Theorem,Inductive,Fixpoint,Record,is},
  sensitive=true,
  morecomment=[s]{(*}{*)},
}
\lstset{
  language=Coq,
  basicstyle=\ttfamily\footnotesize,
  keywordstyle=\bfseries,
  commentstyle=\itshape\color{black!60},
  columns=fullflexible,
  keepspaces=true,
  breaklines=true,
  xleftmargin=1em,
  escapeinside={(@}{@)},
}
% \code{name} typesets an identifier from the source.  It is a url-style command,
% so underscores need no escaping and long names may break after an underscore.
\DeclareUrlCommand\code{\urlstyle{tt}}
\makeatletter
\g@addto@macro\UrlBreaks{\do\_}
\Urlmuskip=0mu\relax  % no letter-spacing in justified columns
\makeatother
\emergencystretch=1.5em

% ---------------------------------------------------------------- macros
\newcommand{\Proc}[2]{\mathsf{Proc}\;#1\;#2}
\newcommand{\NI}{\mathsf{NI}}
\newcommand{\rel}{\mathsf{rel}}
\newcommand{\dis}{\mathsf{dis}}
\newcommand{\todo}[1]{\textcolor{red}{[TODO: #1]}}

% ---------------------------------------------------------------- mechanisation
% Every link to the mechanisation goes through \mech, so the target is set here
% and nowhere else.  \anonymoustrue for submission (Anonymous GitHub mirror of the
% repository), \anonymousfalse for the camera-ready (the repository itself).
% For now the links point to the public repository; switch to \anonymoustrue once
% the mirror exists and its id is filled in below.
\newif\ifanonymous
\anonymousfalse
\ifanonymous
  % TODO: replace XXXX by the id Anonymous GitHub assigns.
  \newcommand{\repobase}{https://anonymous.4open.science/r/NonInterference-XXXX}
  \newcommand{\repofile}[1]{\repobase/#1}
\else
  \newcommand{\repobase}{https://github.com/Tirore96/NonInterference}
  \newcommand{\repofile}[1]{\repobase/blob/main/#1}
\fi
% \m{name}: the Rocq name, linked to the line that defines it.  File and line come
% from mechlinks.tex, which mechlinks.py regenerates from ../theories/*.v before
% every latexmk build, so line numbers stay current.  A name not in the table is
% typeset without a link and reported as a warning.
\newcommand{\mechdef}[3]{%
  \expandafter\def\csname mechfile@#1\endcsname{#2}%
  \expandafter\def\csname mechline@#1\endcsname{#3}}
\DeclareRobustCommand{\m}[1]{%
  \ifcsname mechfile@#1\endcsname
    \href{\repofile{\csname mechfile@#1\endcsname}\#L\csname mechline@#1\endcsname}{\code{#1}}%
  \else
    \code{#1}\PackageWarning{mechlinks}{No definition found for `#1'}%
  \fi}
% \mechfile{path}{text}: link to a whole file of the mechanisation.
\newcommand{\mechfile}[2]{\href{\repofile{#1}}{\code{#2}}}
% Generated by mechlinks.py from ../theories/*.v.  Do not edit.
\mechdef{Clause1}{theories/adequacy.v}{502}
\mechdef{Clause2}{theories/adequacy.v}{506}
\mechdef{Clause3}{theories/adequacy.v}{510}
\mechdef{Clause4}{theories/adequacy.v}{515}
\mechdef{DisTrace}{theories/definitions.v}{827}
\mechdef{HandlerOutput}{theories/definitions.v}{89}
\mechdef{HandlerOutput_hasDecEq}{theories/definitions.v}{100}
\mechdef{HandlerOutput_indDef}{theories/definitions.v}{98}
\mechdef{Interrupt}{theories/definitions.v}{88}
\mechdef{Interrupt_hasDecEq}{theories/definitions.v}{95}
\mechdef{Interrupt_indDef}{theories/definitions.v}{93}
\mechdef{MapTrace}{theories/theorems.v}{47}
\mechdef{NI}{theories/definitions.v}{513}
\mechdef{NI_SNI}{theories/adequacy.v}{419}
\mechdef{NI_ext}{theories/adequacy.v}{167}
\mechdef{NI_l}{theories/definitions.v}{507}
\mechdef{NI_lEquiv}{theories/adequacy.v}{190}
\mechdef{NI_of_SNI}{theories/adequacy.v}{393}
\mechdef{NI_paper}{theories/adequacy.v}{806}
\mechdef{NI_reduceI}{theories/theorems.v}{287}
\mechdef{NI_reduceO}{theories/theorems.v}{306}
\mechdef{ObliviousDis}{theories/adequacy.v}{445}
\mechdef{ObliviousDisF}{theories/adequacy.v}{430}
\mechdef{ObliviousDis_iff}{theories/adequacy.v}{478}
\mechdef{ObliviousDis_oblivious}{theories/adequacy.v}{450}
\mechdef{ObliviousTrace}{theories/definitions.v}{813}
\mechdef{ObliviousTrace_all_outputs_rel}{theories/theorems.v}{815}
\mechdef{Option_presP_all}{theories/definitions.v}{764}
\mechdef{Option_presP_aware}{theories/definitions.v}{771}
\mechdef{Option_presP_nobody}{theories/definitions.v}{778}
\mechdef{Option_presP_nonbot}{theories/definitions.v}{782}
\mechdef{PNI}{theories/adequacy.v}{544}
\mechdef{PNI_of_SNI}{theories/adequacy.v}{715}
\mechdef{Performs}{theories/adequacy.v}{498}
\mechdef{PerformsF}{theories/adequacy.v}{489}
\mechdef{Performs_STrace}{theories/adequacy.v}{549}
\mechdef{Proc}{theories/definitions.v}{153}
\mechdef{Proc_has_input}{theories/theorems.v}{733}
\mechdef{Proc_has_output}{theories/theorems.v}{756}
\mechdef{PublicOutput}{theories/definitions.v}{91}
\mechdef{PublicOutput_hasDecEq}{theories/definitions.v}{110}
\mechdef{PublicOutput_indDef}{theories/definitions.v}{108}
\mechdef{SNI}{theories/adequacy.v}{249}
\mechdef{SNI_l}{theories/adequacy.v}{244}
\mechdef{SNI_of_NI}{theories/adequacy.v}{365}
\mechdef{SNI_of_PNI}{theories/adequacy.v}{793}
\mechdef{SNI_simulation}{theories/adequacy.v}{689}
\mechdef{SRel}{theories/adequacy.v}{600}
\mechdef{SRelF}{theories/adequacy.v}{588}
\mechdef{STrace}{theories/adequacy.v}{241}
\mechdef{STraceF}{theories/adequacy.v}{229}
\mechdef{STrace_follow}{theories/adequacy.v}{332}
\mechdef{STrace_sapp}{theories/adequacy.v}{665}
\mechdef{STrace_scons_out}{theories/adequacy.v}{677}
\mechdef{STrace_simulation}{theories/adequacy.v}{727}
\mechdef{STrace_sprefix}{theories/adequacy.v}{262}
\mechdef{SimulationF}{theories/adequacy.v}{520}
\mechdef{SimulationF_inv}{theories/adequacy.v}{616}
\mechdef{SwiTrace}{theories/theorems.v}{545}
\mechdef{T_in}{theories/models.v}{93}
\mechdef{T_intermediate}{theories/models.v}{234}
\mechdef{T_out}{theories/models.v}{98}
\mechdef{T_out_userview}{theories/models.v}{99}
\mechdef{Trace}{theories/definitions.v}{489}
\mechdef{Trace_STrace}{theories/adequacy.v}{347}
\mechdef{Trace_ext}{theories/adequacy.v}{162}
\mechdef{Trace_imp}{theories/noninterference.v}{65}
\mechdef{Trace_loop_to_p}{theories/theorems.v}{1113}
\mechdef{Trace_map}{theories/models.v}{1045}
\mechdef{Trace_p_to_loop}{theories/theorems.v}{1145}
\mechdef{Trace_stream}{theories/adequacy.v}{358}
\mechdef{Tsum}{theories/models.v}{953}
\mechdef{Tsum'}{theories/models.v}{952}
\mechdef{Ty}{theories/definitions.v}{113}
\mechdef{Ty_eq_dec}{theories/definitions.v}{129}
\mechdef{Ty_hasDecEq}{theories/definitions.v}{124}
\mechdef{Ty_indDef}{theories/definitions.v}{122}
\mechdef{TypeSyscall}{theories/definitions.v}{90}
\mechdef{TypeSyscall_hasDecEq}{theories/definitions.v}{105}
\mechdef{TypeSyscall_indDef}{theories/definitions.v}{103}
\mechdef{all_interrupts}{theories/models.v}{301}
\mechdef{all_outputs_rel}{theories/theorems.v}{808}
\mechdef{all_outputs_rel_insert_inl}{theories/theorems.v}{854}
\mechdef{alternate}{theories/models.v}{893}
\mechdef{aware}{theories/definitions.v}{768}
\mechdef{aware_or_oblivious}{theories/definitions.v}{837}
\mechdef{aware_public}{theories/definitions.v}{1185}
\mechdef{bool_rule_order}{theories/definitions.v}{371}
\mechdef{bool_state}{theories/models.v}{293}
\mechdef{bool_state_equiv}{theories/noninterference.v}{247}
\mechdef{cEquiv}{theories/definitions.v}{76}
\mechdef{cEquiv_of_lEquiv}{theories/adequacy.v}{68}
\mechdef{cEquiv_of_lEquivK}{theories/adequacy.v}{155}
\mechdef{cEquiv_of_lEquiv_dis}{theories/adequacy.v}{130}
\mechdef{cEquiv_of_lEquiv_rel}{theories/adequacy.v}{126}
\mechdef{cEquiv_rule1}{theories/definitions.v}{618}
\mechdef{cEquiv_rule2}{theories/definitions.v}{622}
\mechdef{cEquiv_rule3}{theories/definitions.v}{626}
\mechdef{chc_bits}{theories/models.v}{727}
\mechdef{chc_count}{theories/models.v}{720}
\mechdef{chc_id}{theories/models.v}{723}
\mechdef{check_handler_completed}{theories/models.v}{617}
\mechdef{check_ir_count}{theories/models.v}{621}
\mechdef{check_scheduler}{theories/models.v}{446}
\mechdef{check_scheduler_bits}{theories/models.v}{705}
\mechdef{choice_inv}{theories/models.v}{691}
\mechdef{choice_inv_step}{theories/models.v}{763}
\mechdef{choice_state}{theories/models.v}{686}
\mechdef{cic_close}{theories/models.v}{740}
\mechdef{cic_none}{theories/models.v}{746}
\mechdef{cic_tick}{theories/models.v}{734}
\mechdef{count}{theories/models.v}{285}
\mechdef{count_not_none}{theories/theorems.v}{990}
\mechdef{cur_pid}{theories/models.v}{218}
\mechdef{dI}{theories/models.v}{985}
\mechdef{dI'}{theories/models.v}{955}
\mechdef{dI_o}{theories/models.v}{961}
\mechdef{dI_out}{theories/models.v}{385}
\mechdef{dis_dis_rel}{theories/definitions.v}{590}
\mechdef{dis_eqmaybe_all}{theories/definitions.v}{1166}
\mechdef{dis_eqmaybe_all2}{theories/definitions.v}{1171}
\mechdef{dis_eqmaybe_nobody2}{theories/definitions.v}{1176}
\mechdef{dis_eqpair_L}{theories/definitions.v}{982}
\mechdef{dis_eqpair_L2}{theories/definitions.v}{987}
\mechdef{dis_eqpair_LR}{theories/definitions.v}{992}
\mechdef{dis_eqpair_LR2}{theories/definitions.v}{997}
\mechdef{dis_eqpair_R}{theories/definitions.v}{972}
\mechdef{dis_eqpair_R2}{theories/definitions.v}{977}
\mechdef{dis_eqsum_L}{theories/definitions.v}{1059}
\mechdef{dis_eqsum_L2}{theories/definitions.v}{1063}
\mechdef{dis_mono}{theories/adequacy.v}{50}
\mechdef{dis_rel}{theories/adequacy.v}{54}
\mechdef{dis_rel_dis}{theories/definitions.v}{581}
\mechdef{dis_rel_dis2}{theories/definitions.v}{585}
\mechdef{dis_rel_l}{theories/adequacy.v}{58}
\mechdef{dis_rel_r}{theories/adequacy.v}{62}
\mechdef{dsk_done'}{theories/models.v}{967}
\mechdef{dsk_step'}{theories/models.v}{968}
\mechdef{eqmaybe}{theories/definitions.v}{729}
\mechdef{eqmaybe_all}{theories/definitions.v}{786}
\mechdef{eqmaybe_all_not_bot}{theories/noninterference.v}{183}
\mechdef{eqmaybe_aware}{theories/definitions.v}{794}
\mechdef{eqmaybe_dis}{theories/definitions.v}{728}
\mechdef{eqmaybe_nobody}{theories/definitions.v}{798}
\mechdef{eqmaybe_nonbot}{theories/definitions.v}{790}
\mechdef{eqmaybe_nonbot_at_bot}{theories/noninterference.v}{193}
\mechdef{eqmaybe_nonbot_not_bot}{theories/noninterference.v}{188}
\mechdef{eqmaybe_nonbot_not_bot'}{theories/noninterference.v}{178}
\mechdef{eqpair}{theories/definitions.v}{661}
\mechdef{eqpair_L}{theories/definitions.v}{675}
\mechdef{eqpair_LR}{theories/definitions.v}{683}
\mechdef{eqpair_R}{theories/definitions.v}{679}
\mechdef{eqpair_aux}{theories/definitions.v}{636}
\mechdef{eqsum}{theories/definitions.v}{719}
\mechdef{eqsum_L}{theories/definitions.v}{723}
\mechdef{eqsum_L_llrr}{theories/noninterference.v}{140}
\mechdef{eqsum_L_split}{theories/noninterference.v}{157}
\mechdef{eqsum_aux}{theories/definitions.v}{692}
\mechdef{eqsum_llrr}{theories/noninterference.v}{134}
\mechdef{eqsum_split}{theories/noninterference.v}{146}
\mechdef{f_EP}{theories/definitions.v}{927}
\mechdef{f_EP_comp}{theories/noninterference.v}{307}
\mechdef{f_EP_step_sum}{theories/noninterference.v}{330}
\mechdef{f_NI}{theories/definitions.v}{923}
\mechdef{f_NI_PU}{theories/definitions.v}{925}
\mechdef{f_NI_id}{theories/definitions.v}{1002}
\mechdef{f_NI_snd_eqpair}{theories/definitions.v}{1014}
\mechdef{f_NI_snd_eqpair_LR}{theories/definitions.v}{1019}
\mechdef{f_PU}{theories/definitions.v}{924}
\mechdef{f_PU_id}{theories/definitions.v}{1006}
\mechdef{f_proj}{theories/models.v}{235}
\mechdef{false_equiv}{theories/definitions.v}{516}
\mechdef{false_equiv_aware}{theories/noninterference.v}{232}
\mechdef{false_ic}{theories/models.v}{575}
\mechdef{false_ir_bits}{theories/models.v}{574}
\mechdef{filter_none}{theories/theorems.v}{982}
\mechdef{filter_none_insert_none}{theories/theorems.v}{1001}
\mechdef{filter_none_insert_some}{theories/theorems.v}{1010}
\mechdef{filter_none_trace}{theories/theorems.v}{1026}
\mechdef{filter_none_trace_back}{theories/theorems.v}{1039}
\mechdef{first_ready}{theories/models.v}{479}
\mechdef{first_ready_update_re_sched}{theories/noninterference.v}{406}
\mechdef{follow}{theories/adequacy.v}{329}
\mechdef{fv_NI}{theories/definitions.v}{926}
\mechdef{fv_NI_comp}{theories/noninterference.v}{301}
\mechdef{fv_NI_step_sum}{theories/noninterference.v}{316}
\mechdef{get_bool_state}{theories/models.v}{308}
\mechdef{get_cur_pid}{theories/models.v}{309}
\mechdef{get_ic}{theories/models.v}{314}
\mechdef{get_ic_count}{theories/models.v}{312}
\mechdef{get_ir_bits}{theories/models.v}{322}
\mechdef{get_ir_bits'}{theories/models.v}{316}
\mechdef{get_ir_count}{theories/models.v}{313}
\mechdef{get_ir_mask}{theories/models.v}{328}
\mechdef{get_ir_pending}{theories/models.v}{327}
\mechdef{get_mask'}{theories/models.v}{325}
\mechdef{get_pending'}{theories/models.v}{324}
\mechdef{get_pids}{theories/models.v}{307}
\mechdef{get_prev_pid}{theories/models.v}{310}
\mechdef{get_prev_pid_wrap}{theories/models.v}{493}
\mechdef{get_re_sched}{theories/models.v}{311}
\mechdef{good_no_dI}{theories/models.v}{1056}
\mechdef{good_no_dI_eq}{theories/models.v}{1059}
\mechdef{good_with_dI}{theories/models.v}{1057}
\mechdef{good_with_dI_eq}{theories/models.v}{1062}
\mechdef{handler_completed}{theories/models.v}{602}
\mechdef{handler_completed_reachable}{theories/models.v}{936}
\mechdef{handler_runtime}{theories/models.v}{921}
\mechdef{handler_type}{theories/models.v}{108}
\mechdef{hc_ge2}{theories/models.v}{754}
\mechdef{hc_slice}{theories/models.v}{750}
\mechdef{helper_trace}{theories/noninterference.v}{87}
\mechdef{helper_trace'}{theories/noninterference.v}{73}
\mechdef{high_out}{theories/models.v}{958}
\mechdef{ic}{theories/models.v}{284}
\mechdef{ic_count}{theories/models.v}{292}
\mechdef{ic_equiv}{theories/noninterference.v}{246}
\mechdef{immediate_init}{theories/models.v}{576}
\mechdef{immediate_no_dI'}{theories/models.v}{1001}
\mechdef{immediate_service}{theories/models.v}{578}
\mechdef{immediate_with_dI'}{theories/models.v}{1002}
\mechdef{in_equiv}{theories/noninterference.v}{31}
\mechdef{in_equiv_eq}{theories/noninterference.v}{269}
\mechdef{initial_pid}{theories/models.v}{219}
\mechdef{initiate_handler}{theories/models.v}{472}
\mechdef{initiate_ir}{theories/models.v}{610}
\mechdef{initiate_ir_count}{theories/models.v}{715}
\mechdef{initiate_ir_mask}{theories/models.v}{711}
\mechdef{initiate_next}{theories/models.v}{496}
\mechdef{initiate_prev_pid}{theories/models.v}{494}
\mechdef{initiate_scheduler}{theories/models.v}{491}
\mechdef{inj_pair2_Ty}{theories/definitions.v}{132}
\mechdef{inr_or_def}{theories/models.v}{254}
\mechdef{insert}{theories/definitions.v}{495}
\mechdef{insert_out}{theories/theorems.v}{1020}
\mechdef{insert_oversize}{theories/adequacy.v}{320}
\mechdef{interp}{theories/definitions.v}{136}
\mechdef{inv_loopI}{theories/definitions.v}{308}
\mechdef{inv_loopO}{theories/definitions.v}{358}
\mechdef{inv_mapI}{theories/definitions.v}{285}
\mechdef{inv_mapO}{theories/definitions.v}{315}
\mechdef{inv_maybeI_None}{theories/definitions.v}{297}
\mechdef{inv_maybeI_Some}{theories/definitions.v}{300}
\mechdef{inv_maybeO}{theories/definitions.v}{342}
\mechdef{inv_outI}{theories/definitions.v}{282}
\mechdef{inv_outO}{theories/definitions.v}{312}
\mechdef{inv_parI}{theories/definitions.v}{304}
\mechdef{inv_parO}{theories/definitions.v}{349}
\mechdef{inv_staI}{theories/definitions.v}{289}
\mechdef{inv_staO}{theories/definitions.v}{322}
\mechdef{inv_swiI}{theories/definitions.v}{293}
\mechdef{inv_swiO_false}{theories/definitions.v}{330}
\mechdef{inv_swiO_true}{theories/definitions.v}{334}
\mechdef{ir_bits}{theories/models.v}{282}
\mechdef{ir_count}{theories/models.v}{283}
\mechdef{ir_dis}{theories/noninterference.v}{30}
\mechdef{ir_handler}{theories/models.v}{121}
\mechdef{ir_handler_pid}{theories/models.v}{455}
\mechdef{ir_ready}{theories/models.v}{476}
\mechdef{irq}{theories/models.v}{297}
\mechdef{is_handler_pid}{theories/models.v}{482}
\mechdef{is_inl}{theories/definitions.v}{1181}
\mechdef{is_inr}{theories/definitions.v}{1182}
\mechdef{is_ir_out_done}{theories/models.v}{395}
\mechdef{is_sched_out}{theories/models.v}{441}
\mechdef{lEquiv}{theories/adequacy.v}{30}
\mechdef{lEquiv_of_cEquiv}{theories/adequacy.v}{113}
\mechdef{lEquiv_of_cEquivK}{theories/adequacy.v}{144}
\mechdef{lEquiv_of_cEquiv_obs}{theories/adequacy.v}{135}
\mechdef{lEquiv_of_cEquiv_unobs}{theories/adequacy.v}{139}
\mechdef{level}{theories/definitions.v}{21}
\mechdef{levelPred}{theories/definitions.v}{599}
\mechdef{level_not}{theories/definitions.v}{1087}
\mechdef{lobs}{theories/adequacy.v}{40}
\mechdef{loop_NI}{theories/theorems.v}{1216}
\mechdef{loop_trace}{theories/theorems.v}{1108}
\mechdef{low_out}{theories/models.v}{957}
\mechdef{lrel_of_cEquiv}{theories/adequacy.v}{85}
\mechdef{lrel_of_cEquiv_refl}{theories/adequacy.v}{93}
\mechdef{lrel_of_cEquiv_sym}{theories/adequacy.v}{97}
\mechdef{lrel_of_cEquiv_trans}{theories/adequacy.v}{103}
\mechdef{lt_ins_inl}{theories/theorems.v}{1168}
\mechdef{lt_ins_inl_inv}{theories/theorems.v}{1186}
\mechdef{lthread}{theories/theorems.v}{369}
\mechdef{lthread_insert_dis}{theories/theorems.v}{459}
\mechdef{lthread_remove_dis}{theories/theorems.v}{478}
\mechdef{lthread_stable}{theories/theorems.v}{376}
\mechdef{lthread_swap}{theories/theorems.v}{439}
\mechdef{lunobs}{theories/adequacy.v}{42}
\mechdef{map_NI}{theories/theorems.v}{129}
\mechdef{map_insert}{theories/theorems.v}{1247}
\mechdef{map_tr}{theories/models.v}{1044}
\mechdef{map_trace}{theories/theorems.v}{52}
\mechdef{map_trace_insert}{theories/theorems.v}{69}
\mechdef{map_trace_insert2}{theories/theorems.v}{84}
\mechdef{map_trace_insert3}{theories/theorems.v}{99}
\mechdef{map_trace_insert4}{theories/theorems.v}{122}
\mechdef{map_trace_remove3}{theories/theorems.v}{111}
\mechdef{mask}{theories/models.v}{280}
\mechdef{mask_most}{theories/models.v}{592}
\mechdef{masks_set}{theories/models.v}{370}
\mechdef{masks_setE}{theories/models.v}{700}
\mechdef{maybe_NI}{theories/theorems.v}{1059}
\mechdef{model}{theories/models.v}{527}
\mechdef{model_immediate}{theories/models.v}{583}
\mechdef{model_immediate_concrete}{theories/models.v}{927}
\mechdef{model_immediate_not_NI}{theories/noninterference.v}{102}
\mechdef{model_immediate_not_NI_paper}{theories/adequacy.v}{832}
\mechdef{model_sliced}{theories/models.v}{651}
\mechdef{model_sliced_NI}{theories/noninterference.v}{550}
\mechdef{model_sliced_NI_paper}{theories/adequacy.v}{819}
\mechdef{model_sliced_concrete}{theories/models.v}{928}
\mechdef{model_sliced_concrete_NI}{theories/noninterference.v}{1732}
\mechdef{model_sliced_userview}{theories/models.v}{864}
\mechdef{model_sliced_userview_NI}{theories/noninterference.v}{1707}
\mechdef{model_sliced_userview_concrete}{theories/models.v}{929}
\mechdef{model_sliced_userview_concrete_NI}{theories/noninterference.v}{1738}
\mechdef{monotone_ObliviousDisF}{theories/adequacy.v}{436}
\mechdef{monotone_PerformsF}{theories/adequacy.v}{494}
\mechdef{monotone_SRelF}{theories/adequacy.v}{595}
\mechdef{monotone_STraceF}{theories/adequacy.v}{236}
\mechdef{monotone_SimulationF}{theories/adequacy.v}{527}
\mechdef{no_output}{theories/models.v}{401}
\mechdef{nop_done'}{theories/models.v}{978}
\mechdef{nop_ir_out}{theories/models.v}{390}
\mechdef{nop_o}{theories/models.v}{960}
\mechdef{nop_step'}{theories/models.v}{977}
\mechdef{not_booleq}{theories/definitions.v}{553}
\mechdef{oblivious}{theories/definitions.v}{834}
\mechdef{oblivious_NI_l}{theories/theorems.v}{872}
\mechdef{oblivious_ObliviousDis}{theories/adequacy.v}{464}
\mechdef{oblivious_all_dis}{theories/theorems.v}{973}
\mechdef{oblivious_at}{theories/definitions.v}{818}
\mechdef{oblivious_at_reduceI}{theories/theorems.v}{784}
\mechdef{oblivious_at_reduceO}{theories/theorems.v}{792}
\mechdef{oblivious_at_reduceO_rel}{theories/theorems.v}{800}
\mechdef{oblivious_at_trace_any}{theories/theorems.v}{821}
\mechdef{oblivious_at_trace_rel}{theories/theorems.v}{845}
\mechdef{oblivious_class}{theories/definitions.v}{823}
\mechdef{option_rule_order}{theories/definitions.v}{367}
\mechdef{or_bool_state}{theories/models.v}{359}
\mechdef{or_ic}{theories/models.v}{354}
\mechdef{or_ir_bits}{theories/models.v}{353}
\mechdef{order}{theories/definitions.v}{67}
\mechdef{order_bot}{theories/definitions.v}{545}
\mechdef{out_NI}{theories/theorems.v}{17}
\mechdef{out_equiv}{theories/noninterference.v}{54}
\mechdef{out_equiv_not_bot}{theories/noninterference.v}{208}
\mechdef{out_equiv_userview}{theories/models.v}{870}
\mechdef{out_get}{theories/models.v}{980}
\mechdef{out_get'}{theories/models.v}{970}
\mechdef{out_nop}{theories/models.v}{981}
\mechdef{out_nop'}{theories/models.v}{971}
\mechdef{out_syscall}{theories/models.v}{982}
\mechdef{out_syscall'}{theories/models.v}{972}
\mechdef{p_priv_concrete}{theories/models.v}{906}
\mechdef{p_priv_concrete_NI}{theories/noninterference.v}{427}
\mechdef{p_pub_concrete}{theories/models.v}{884}
\mechdef{p_pub_concrete_NI}{theories/noninterference.v}{420}
\mechdef{pair_rewr}{theories/definitions.v}{1080}
\mechdef{par_NI}{theories/theorems.v}{1297}
\mechdef{par_to_proj}{theories/theorems.v}{1251}
\mechdef{parse_output}{theories/models.v}{409}
\mechdef{pending}{theories/models.v}{281}
\mechdef{pids}{theories/models.v}{288}
\mechdef{pids_equiv}{theories/noninterference.v}{243}
\mechdef{pl1}{theories/theorems.v}{1242}
\mechdef{pl2}{theories/theorems.v}{1244}
\mechdef{pool}{theories/models.v}{244}
\mechdef{pool_concrete}{theories/models.v}{925}
\mechdef{pool_input}{theories/models.v}{404}
\mechdef{presP}{theories/definitions.v}{600}
\mechdef{prev_pid}{theories/models.v}{287}
\mechdef{private_bot}{theories/noninterference.v}{198}
\mechdef{private_equiv}{theories/definitions.v}{561}
\mechdef{private_not_bot}{theories/noninterference.v}{168}
\mechdef{private_not_bot'}{theories/noninterference.v}{173}
\mechdef{private_pending_equiv}{theories/noninterference.v}{244}
\mechdef{process_pool}{theories/models.v}{192}
\mechdef{projI}{theories/theorems.v}{199}
\mechdef{projIl}{theories/theorems.v}{212}
\mechdef{projIl_insert}{theories/theorems.v}{359}
\mechdef{projIl_insert_inv}{theories/theorems.v}{326}
\mechdef{projO}{theories/theorems.v}{205}
\mechdef{projOl}{theories/theorems.v}{218}
\mechdef{projOl_insert}{theories/theorems.v}{352}
\mechdef{proj_to_par}{theories/theorems.v}{1283}
\mechdef{public_bits_equiv}{theories/noninterference.v}{245}
\mechdef{public_equiv}{theories/definitions.v}{530}
\mechdef{public_equiv_eq}{theories/noninterference.v}{205}
\mechdef{re_sched}{theories/models.v}{286}
\mechdef{reactive_system}{theories/models.v}{261}
\mechdef{ready_aux}{theories/noninterference.v}{339}
\mechdef{ready_aux2}{theories/noninterference.v}{361}
\mechdef{ready_auxP}{theories/noninterference.v}{341}
\mechdef{ready_cases}{theories/noninterference.v}{359}
\mechdef{ready_cases2}{theories/noninterference.v}{365}
\mechdef{ready_cases2P}{theories/noninterference.v}{374}
\mechdef{ready_casesP}{theories/noninterference.v}{369}
\mechdef{reduceI}{theories/definitions.v}{177}
\mechdef{reduceI_det}{theories/theorems.v}{278}
\mechdef{reduceI_fun}{theories/definitions.v}{253}
\mechdef{reduceI_step}{theories/definitions.v}{231}
\mechdef{reduceO}{theories/definitions.v}{189}
\mechdef{reduceO_det}{theories/theorems.v}{281}
\mechdef{reduceO_fun}{theories/definitions.v}{263}
\mechdef{reduceO_step}{theories/definitions.v}{238}
\mechdef{rel_eq}{theories/definitions.v}{802}
\mechdef{rel_eqmaybe_all}{theories/definitions.v}{1082}
\mechdef{rel_eqmaybe_all2}{theories/definitions.v}{1148}
\mechdef{rel_eqmaybe_all_aux}{theories/definitions.v}{1135}
\mechdef{rel_eqmaybe_aware2}{theories/definitions.v}{1157}
\mechdef{rel_eqmaybe_nobody2}{theories/definitions.v}{1130}
\mechdef{rel_eqmaybe_nonbot}{theories/definitions.v}{1117}
\mechdef{rel_eqmaybe_nonbot2}{theories/definitions.v}{1126}
\mechdef{rel_eqmaybe_nonbot_aux}{theories/definitions.v}{1104}
\mechdef{rel_eqpair}{theories/definitions.v}{931}
\mechdef{rel_eqpair2}{theories/definitions.v}{936}
\mechdef{rel_eqpair_L2}{theories/definitions.v}{952}
\mechdef{rel_eqpair_L2'}{theories/definitions.v}{957}
\mechdef{rel_eqpair_LR}{theories/definitions.v}{962}
\mechdef{rel_eqpair_LR2}{theories/definitions.v}{967}
\mechdef{rel_eqpair_R2}{theories/definitions.v}{942}
\mechdef{rel_eqpair_R2'}{theories/definitions.v}{947}
\mechdef{rel_eqpair_to_L}{theories/definitions.v}{1068}
\mechdef{rel_eqsum}{theories/definitions.v}{1027}
\mechdef{rel_eqsum'}{theories/definitions.v}{1035}
\mechdef{rel_eqsum2}{theories/definitions.v}{1031}
\mechdef{rel_eqsum2'}{theories/definitions.v}{1039}
\mechdef{rel_eqsum_L}{theories/definitions.v}{1043}
\mechdef{rel_eqsum_L'}{theories/definitions.v}{1051}
\mechdef{rel_eqsum_L2}{theories/definitions.v}{1047}
\mechdef{rel_eqsum_L2'}{theories/definitions.v}{1055}
\mechdef{rel_mono}{theories/adequacy.v}{46}
\mechdef{rel_refl}{theories/definitions.v}{602}
\mechdef{rel_sym}{theories/definitions.v}{606}
\mechdef{rel_trans}{theories/definitions.v}{612}
\mechdef{remove}{theories/definitions.v}{501}
\mechdef{restore_invariant}{theories/models.v}{640}
\mechdef{restore_invariant_id}{theories/models.v}{808}
\mechdef{restore_invariant_noop}{theories/models.v}{854}
\mechdef{runS}{theories/adequacy.v}{573}
\mechdef{runS_Performs}{theories/adequacy.v}{579}
\mechdef{runS_SRel}{theories/adequacy.v}{603}
\mechdef{runT}{theories/adequacy.v}{340}
\mechdef{runT_cat}{theories/adequacy.v}{661}
\mechdef{sans_timer}{theories/models.v}{303}
\mechdef{sapp}{theories/adequacy.v}{215}
\mechdef{sapp_cat}{theories/adequacy.v}{653}
\mechdef{sapp_scons}{theories/adequacy.v}{657}
\mechdef{save_cur_to_prev}{theories/models.v}{467}
\mechdef{sched_o}{theories/models.v}{959}
\mechdef{sched_priv'}{theories/models.v}{975}
\mechdef{sched_pub'}{theories/models.v}{974}
\mechdef{scheduler}{theories/models.v}{911}
\mechdef{scheduler_NI}{theories/noninterference.v}{464}
\mechdef{scheduler_pid}{theories/models.v}{489}
\mechdef{seqtype}{theories/models.v}{988}
\mechdef{seqtype'}{theories/models.v}{987}
\mechdef{set_masks}{theories/models.v}{361}
\mechdef{set_otherIs}{theories/models.v}{369}
\mechdef{set_pending}{theories/models.v}{438}
\mechdef{set_tI}{theories/models.v}{368}
\mechdef{shd}{theories/adequacy.v}{201}
\mechdef{simp_pair1}{theories/definitions.v}{1073}
\mechdef{simp_pair2}{theories/definitions.v}{1076}
\mechdef{simulation}{theories/adequacy.v}{540}
\mechdef{simulation_SRel}{theories/adequacy.v}{622}
\mechdef{simulation_STrace}{theories/adequacy.v}{559}
\mechdef{simulation_del_n}{theories/adequacy.v}{775}
\mechdef{simulation_ins_n}{theories/adequacy.v}{756}
\mechdef{simulation_sub_n}{theories/adequacy.v}{737}
\mechdef{sinsert}{theories/adequacy.v}{221}
\mechdef{sinsert_sapp}{theories/adequacy.v}{312}
\mechdef{slice_runs}{theories/models.v}{922}
\mechdef{sliced_init}{theories/models.v}{594}
\mechdef{sliced_inv}{theories/models.v}{680}
\mechdef{sliced_inv_step}{theories/models.v}{815}
\mechdef{sliced_no_dI'}{theories/models.v}{1003}
\mechdef{sliced_run}{theories/models.v}{844}
\mechdef{sliced_run_inv}{theories/models.v}{848}
\mechdef{sliced_service}{theories/models.v}{646}
\mechdef{sliced_with_dI'}{theories/models.v}{1004}
\mechdef{slot_I}{theories/models.v}{131}
\mechdef{slot_O}{theories/models.v}{142}
\mechdef{slot_initial}{theories/models.v}{232}
\mechdef{slot_pid}{theories/models.v}{221}
\mechdef{slot_pid_ir_handler_pid}{theories/models.v}{463}
\mechdef{slot_procs}{theories/models.v}{163}
\mechdef{slot_procs_concrete}{theories/models.v}{924}
\mechdef{sprefix}{theories/adequacy.v}{209}
\mechdef{sprefix_STrace}{theories/adequacy.v}{273}
\mechdef{sprefix_sapp}{theories/adequacy.v}{309}
\mechdef{sprefix_sinsert_gt}{theories/adequacy.v}{299}
\mechdef{sprefix_sinsert_le}{theories/adequacy.v}{292}
\mechdef{sta_NI}{theories/theorems.v}{499}
\mechdef{sta_comp}{theories/noninterference.v}{274}
\mechdef{sta_conv}{theories/theorems.v}{391}
\mechdef{sta_proj}{theories/theorems.v}{226}
\mechdef{sta_proj_insert}{theories/theorems.v}{245}
\mechdef{sta_proj_lthread}{theories/theorems.v}{414}
\mechdef{state_step}{theories/models.v}{509}
\mechdef{state_type}{theories/models.v}{305}
\mechdef{state_type_equiv}{theories/noninterference.v}{248}
\mechdef{state_type_equiv_not_bot}{theories/noninterference.v}{251}
\mechdef{stepI}{theories/definitions.v}{202}
\mechdef{stepO}{theories/definitions.v}{216}
\mechdef{step_sum}{theories/models.v}{425}
\mechdef{stl}{theories/adequacy.v}{202}
\mechdef{stream}{theories/adequacy.v}{198}
\mechdef{sunfold}{theories/adequacy.v}{204}
\mechdef{sunfoldE}{theories/adequacy.v}{206}
\mechdef{swi_NI}{theories/theorems.v}{954}
\mechdef{swi_NI_class}{theories/theorems.v}{904}
\mechdef{swi_conv}{theories/theorems.v}{701}
\mechdef{swi_trace}{theories/theorems.v}{551}
\mechdef{swi_trace_insert}{theories/theorems.v}{581}
\mechdef{swi_trace_insert_conv}{theories/theorems.v}{607}
\mechdef{swi_trace_remove}{theories/theorems.v}{668}
\mechdef{swi_trace_swap}{theories/theorems.v}{640}
\mechdef{tI}{theories/models.v}{984}
\mechdef{tI'}{theories/models.v}{956}
\mechdef{tI_o}{theories/models.v}{962}
\mechdef{tI_out}{theories/models.v}{380}
\mechdef{threaded}{theories/theorems.v}{252}
\mechdef{time_slice}{theories/models.v}{600}
\mechdef{times_n}{theories/models.v}{84}
\mechdef{times_on}{theories/models.v}{90}
\mechdef{timeslice_live}{theories/models.v}{628}
\mechdef{tmr_done'}{theories/models.v}{964}
\mechdef{tmr_step'}{theories/models.v}{965}
\mechdef{trace_immediate_no_dI'}{theories/models.v}{1029}
\mechdef{trace_immediate_with_dI'}{theories/models.v}{1032}
\mechdef{trace_no_dI}{theories/models.v}{1065}
\mechdef{trace_sliced_no_dI'}{theories/models.v}{1035}
\mechdef{trace_sliced_with_dI'}{theories/models.v}{1038}
\mechdef{trace_with_dI}{theories/models.v}{1072}
\mechdef{unit_proc}{theories/models.v}{154}
\mechdef{unset_masks}{theories/models.v}{363}
\mechdef{unset_masks_sans}{theories/models.v}{365}
\mechdef{unset_tI}{theories/models.v}{367}
\mechdef{update_bool_state}{theories/models.v}{331}
\mechdef{update_cur_pid}{theories/models.v}{332}
\mechdef{update_ic}{theories/models.v}{345}
\mechdef{update_ic_count}{theories/models.v}{343}
\mechdef{update_ir_bits}{theories/models.v}{346}
\mechdef{update_ir_bits'}{theories/models.v}{336}
\mechdef{update_ir_count}{theories/models.v}{344}
\mechdef{update_ir_mask}{theories/models.v}{350}
\mechdef{update_ir_pending}{theories/models.v}{348}
\mechdef{update_pending_rel}{theories/noninterference.v}{513}
\mechdef{update_pids}{theories/models.v}{330}
\mechdef{update_prev_pid}{theories/models.v}{333}
\mechdef{update_re_sched}{theories/models.v}{334}
\mechdef{update_re_sched_rel}{theories/noninterference.v}{534}
\mechdef{w_None}{theories/models.v}{983}
\mechdef{xor}{theories/definitions.v}{163}
\mechdef{xorC}{theories/definitions.v}{168}
\mechdef{xorFalse}{theories/definitions.v}{171}
\mechdef{xorK}{theories/definitions.v}{165}
\mechdef{xorTrue}{theories/definitions.v}{174}
\mechdef{zerol}{theories/definitions.v}{840}


% IEEEtran defines no theorem environments.
\newtheorem{theorem}{Theorem}
\newtheorem{corollary}{Corollary}
\theoremstyle{definition}
\newtheorem{definition}{Definition}

\begin{document}

\title{Timing-Sensitive Non-Interference\\for Interrupt Handling,\\
  Compositionally and Mechanically}

% TODO: complete the author list and affiliations.
\author{\IEEEauthorblockN{Willard Rafnsson}
\IEEEauthorblockA{IT University of Copenhagen, Denmark\\
wilr@itu.dk}
\and
\IEEEauthorblockN{Dawit Tirore}
\IEEEauthorblockA{IT University of Copenhagen, Denmark\\
dati@itu.dk}}

\maketitle

\begin{abstract}
The security of a computer system rests on its operating system, which ties
software to hardware and decides which process runs, when, and on whose behalf. A key part of
this is interrupt handling. Servicing an interrupt takes time away from whichever
process was scheduled. If a public process keeps issuing requests whose results an
attacker can see, the handler shows up as a gap in that stream, and the attacker
learns that the interrupt occurred. Such timing leaks through
interrupt handling arise in several operating systems
\todo{examples, with citations}. Leaks of this kind are studied in information flow
control, where their absence is stated as non-interference. In this paper we use
the process calculus of Rafnsson et al.\ to model interrupt handling as a process,
and prove in Rocq that it is interfering: a secret disk interrupt is observable
through the schedule alone. We then model the corresponding non-interfering
variant, obtained by altering how interrupts are processed: they are serviced only
inside publicly timed slices, and unused slice time is filled with a dummy handler.
We prove it timing-sensitive non-interfering for any scheduler and user processes
that are themselves non-interfering. The implementation of the modified interrupt
processing is carried out in separate work \todo{cite X}.
\end{abstract}

\begin{IEEEkeywords}
non-interference, timing channels, interrupts, process calculi, compositional
verification, Rocq
\end{IEEEkeywords}

\section{Introduction}\label{sec:intro}

The security of a computer system rests on its operating system. The operating
system ties software to hardware: it owns the devices, decides which process runs and when, and stands
between processes that must not learn about each other. Whatever an operating
system lets one process observe about another is a potential leak, and that
includes \emph{when} things happen, not only what values are computed.

A key part of this mediation is interrupt handling. Interrupts are how an operating
system learns that something happened: a disk transfer completed, a packet arrived,
a timer expired. Some of these events are secret. Whether a particular process just
received data from the disk may reveal what it is doing, so an operating system
that keeps such an event secret has to hide that it happened at all.

Ordinary interrupt handling does not hide it. As soon as the interrupt can be
serviced, its handler runs, displacing whatever process was scheduled; the handler
takes a few steps, and then the displaced process resumes. Suppose a public process keeps
issuing requests whose results an attacker can see. While the handler runs, those
requests stop, and the attacker sees a gap in the stream exactly as long as the
handler (\cref{fig:leak}). The interrupt has leaked through the \emph{schedule}.
Timing leaks through interrupt handling of this kind arise in several operating
systems \todo{examples, with citations}.

This is the example we follow throughout the paper. The secret is \emph{whether a
disk interrupt occurred}. A public process issues requests whose results the
attacker sees, and a private process, the legitimate consumer of the disk event,
reacts to it. The attacker may learn nothing about the disk interrupt. \Cref{sec:model,sec:designs,sec:proof}
model this system, exhibit the leak, and remove it.

\begin{figure}[t]
\centering
\begin{tabular}{lcccc}
\toprule
step & 1 & 2 & 3 & 4 \\
\midrule
no disk interrupt   & \code{Get} & \code{Get} & \code{Get} & \code{Get} \\
disk interrupt before step 1 & \code{Get} & ---   & ---   & \code{Get} \\
\bottomrule
\end{tabular}
\caption{What the attacker sees of the public process under ordinary interrupt
  handling, with handlers of two steps. The two runs differ only in whether a disk
  interrupt arrives, which the attacker may not learn. The interrupt sets a pending
  bit; step 1 was already scheduled, so the disk handler runs at steps 2 and 3, and
  the public process does not run then. The mechanised counterexample (\cref{thm:leak}) uses steps 1 and 2.}
\label{fig:leak}
\end{figure}

Leaks of this kind are the subject of information flow
control~\cite{sabelfeld2003}, where the absence of a leak is stated as
non-interference~\cite{goguen1982}: varying what the attacker may not see must not
change what it does see. For timing leaks, ``what it sees'' must include when each
output appears~\cite{agat2000}.

In this paper we use the process calculus of Rafnsson, Jia and
Bauer~\cite{rafnsson2017} to state and prove such properties of interrupt handling.
Its notion of non-interference is timing-sensitive, and its non-interference
theorems compose: if the parts of a process are non-interfering at suitable
equivalences, so is the whole. We model an operating system with interrupt handling
as one closed process in this calculus. Inside it, the two user processes (one
public, one private), the scheduler and one handler per kind of interrupt run as
concurrent processes in a pool, and a global state, holding the interrupt
controller, decides at each step which of them runs. Interrupts arrive as inputs to
the whole system, and at each step the pool emits the output of the process that
ran. We prove in
Rocq\footnote{The mechanisation is available at \url{\repobase}.} that ordinary
interrupt handling is interfering (\m{model_immediate_not_NI}): the leak of
\cref{fig:leak} is a counterexample.

We then model the corresponding non-interfering variant, obtained by altering how
interrupts are processed. The idea is to take away a secret interrupt's power to
choose \emph{when} its handler runs. Interrupts are serviced only inside time slices
opened by the public timer, every handler run takes the same number of steps, and
any part of a slice that no real interrupt claims is filled with a dummy handler,
the NOP handler. A secret handler run then replaces a NOP run instead of displacing a user
process, and the public schedule is the same whether or not the interrupt occurred.
We prove this variant non-interfering (\m{model_sliced_NI}), even to an attacker who
sees, at every step, which process in the pool ran, the scheduler and the handlers
included. Both variants are instances of a single generic model and differ only
in their \emph{interrupt policy}: the initial state and how handler completion is
serviced. The result is general in two ways. It holds for every scheduler
and every pair of user processes that are themselves non-interfering. And it holds
for every choice of the two numbers that fix the timing of the policy: the number
of steps one handler run takes, and the number of handler runs that fit in one
time slice.

The model fixes \emph{what} the modified interrupt processing must do; how to
implement it in an operating system is carried out in separate work
\todo{cite X}.

\paragraph{Contributions}
\begin{itemize}
\item A mechanised, timing-sensitive non-interference proof for interrupt handling
  (\cref{sec:designs,sec:proof}), and a mechanised counterexample for ordinary
  interrupt handling (\cref{sec:leak}). The proof is compositional: it applies one generic
  theorem per constructor of the model, and leaves one substantial obligation
  specific to interrupt handling (\cref{sec:fvni}).
\item Requirements, forced by the proof, on any implementation of interrupt
  handling that is to be non-interfering (\cref{sec:proof}). The arrival of a
  secret interrupt may only record that it is pending; every decision it influences, such as starting its handler, must be
  deferred to a point fixed by public time. Whether a secret handler ran must be
  indistinguishable from whether a dummy handler ran, so the time reserved for
  handlers must be fixed in advance and always used up. And the interrupt
  controller must be split by secrecy: which interrupts are pending is secret,
  while the masks that decide what runs next may depend only on public
  information. The implementation in \todo{cite X} follows these requirements.
\item A mechanisation of the calculus and its composition theorems. Our
  definitions depart from the original in four places, and for each we prove that
  the departure changes nothing or only weakens a side condition
  (\cref{sec:mechanisation}).
\end{itemize}

\Cref{sec:background} recalls the calculus and the definition of non-interference.
\Cref{sec:model} presents the generic model and \cref{sec:designs} the two interrupt policies.
\Cref{sec:proof} gives the proof, \cref{sec:generality} states precisely how general
it is, and \cref{sec:mechanisation} describes the mechanisation.
\Cref{sec:related,sec:conclusion} discuss related work and limitations.

\section{Background: Composable Timing-Sensitive Non-Interference}\label{sec:background}

This section recalls the calculus and definitions of Rafnsson et
al.~\cite{rafnsson2017}, in the form we mechanise. The form differs from the
original in four places, discussed in \cref{sec:departures}.

\subsection{Processes, traces and time}

A process $p : \Proc{I}{O}$ has an input type $I$ and an output type $O$. It moves
in one of two ways, by receiving an input ($p \xrightarrow{?i} p'$) or by emitting
an output ($p \xrightarrow{!o} p'$), in any order. Both kinds of step are total and
deterministic: in every state a process can receive any input and can emit, and the
successor (and on an output step the value emitted) is determined. A \emph{trace} is
the sequence of labels along a finite run.

Time is counted in output steps. Every output step is one tick of a clock the
observer shares, and on every tick a process emits a value; there is no silent tick.
Where nothing happens, the value says so (in our models a component that did not run
emits \code{None}), and that absence is itself observable. Input steps take no time. This
is what makes the setting timing-sensitive: delaying an output by one step changes
the trace exactly as changing its value would.

Processes are built from seven constructors (\cref{tab:constructors}). The
operating system of \cref{sec:model} is a single closed term over them.

\begin{table}[t]
\centering
\caption{The constructors, by their behaviour on an input and on an output step.}
\label{tab:constructors}
\small
\begin{tabularx}{\columnwidth}{@{}lL@{}}
\toprule
constructor & behaviour \\
\midrule
\code{out o} & receives anything and is unchanged; emits the constant \code{o} \\
\code{map f g p} & hands \code{p} the input \code{f i}; emits \code{g o} when \code{p} emits \code{o} \\
\code{sta f g v p} & a state cell \code{v}: on input \code{i} the cell becomes \code{f i v} and \code{p}
  receives the pair; on output \code{o} of \code{p} the cell becomes \code{g o v} and the pair
  of new cell and \code{o} is emitted \\
\code{swi b p} & a one-bit switch: in phase \code{false} it emits \code{None} and \code{p} does not
  step; in phase \code{true} it emits \code{Some o} for \code{p}'s tagged output \code{(c,o)}, and
  the tag \code{c} and the inputs flip the phase \\
\code{par p1 p2} & both receive every input; both emit, and the pair is emitted \\
\code{loop p} & feedback: \code{p} emits \code{o}, \code{o} is fed back to \code{p} as an input, and \code{o}
  is emitted \\
\code{maybe p} & optional input: \code{p} ignores \code{None} and receives \code{i} on \code{Some i} \\
\bottomrule
\end{tabularx}
\end{table}

\subsection{Characterised equivalences and non-interference}

Security is stated relative to a lattice of levels. Everything in this paper can be
read with two: the attacker $\bot$, and $\top$, who sees everything
(\cref{sec:generality} discusses arbitrary lattices). Each input and output type
carries a \emph{characterised equivalence} (\m{cEquiv}), a record of two
level-indexed predicates:
\begin{itemize}
\item $\rel\;l\;x\;y$: an observer at $l$ cannot tell $x$ from $y$;
\item $\dis\;l\;x$: an observer at $l$ does not see $x$ at all, so $x$ may be
  inserted or removed freely.
\end{itemize}
$\rel\;l$ is an equivalence, and $\dis\;l$ is \emph{characterised} as one of its
classes:
\[
\dis\;l\;a_0 \;\Rightarrow\; \forall a_1.\;
  (\dis\;l\;a_1 \Leftrightarrow \rel\;l\;a_0\;a_1).
\]
So at each level the type is partitioned into classes the observer cannot tell
apart, and one of them, the \emph{distinguished class}, holds the values the
observer does not see. It is empty when the level hides nothing.

Let $\mathit{insert}\;n\;i\;t$ place input $i$ at position $n$ of trace $t$, and let
$\mathsf{Trace}$ compare outputs up to $\rel_{O}\;l$.

\begin{definition}[Non-interference]
$\NI\;\mathit{IRel}\;\mathit{ORel}\;p$ holds when for every level $l$, every trace
$t$, position $n$ and inputs $i, i'$:
\begin{enumerate}
\item (substitution) if $\rel_I\;l\;i\;i'$ and $\mathit{insert}\;n\;i\;t$ is a trace
  of $p$, then so is $\mathit{insert}\;n\;i'\;t$;
\item (insertion) if $\dis_I\;l\;i$ and $t$ is a trace of $p$, then so is
  $\mathit{insert}\;n\;i\;t$;
\item (deletion) if $\dis_I\;l\;i$ and $\mathit{insert}\;n\;i\;t$ is a trace of
  $p$, then so is $t$.
\end{enumerate}
\end{definition}

The set of traces visible at $l$ is thus unchanged by anything the observer at $l$
may not see. All the content is at $\bot$: at $\top$ nothing is unobservable and
indistinguishable means equal, so the clauses hold trivially.

\subsection{Equivalences}

Two equivalences do most of the work. Under \m{public_equiv} nothing is unobservable
and values are related iff equal, at every level. Under \m{private_equiv} every value
is unobservable at $\bot$, where all values are related, and the equivalence is
\m{public_equiv} at $\top$.

Composite equivalences are built from these. \m{eqpair} relates pairs componentwise,
and its variants \m{eqpair_L}, \m{eqpair_R} and \m{eqpair_LR} make the pair secret when
its left, right, or both components are. \m{eqsum} relates a sum tag by tag and never
relates an \code{inl} to an \code{inr}, so the tag stays public even when the payload is not;
\m{eqsum_L} makes an \code{inl} secret exactly when its payload is. The \m{eqmaybe} family
relates \code{Option} values, deferring to an inner equivalence on \code{Some} and differing
in who can see \code{None} (\cref{tab:eqmaybe}). By the characterisation law, \code{Some v}
and \code{None} are related at $l$ exactly when both are unobservable there.

\begin{table}[t]
\centering
\caption{Who can see \texttt{None} under the \texttt{eqmaybe} family. To the
  attacker, \texttt{eqmaybe\_private} and \texttt{eqmaybe\_hidden} agree.}
\label{tab:eqmaybe}
\small
\begin{tabularx}{\columnwidth}{@{}lLL@{}}
\toprule
equivalence & $\bot$ sees \texttt{None} & $\top$ sees \texttt{None} \\
\midrule
\m{eqmaybe}         & yes & yes \\
\m{eqmaybe_private} & no  & yes \\
\m{eqmaybe_hidden}  & no  & no \\
\code{eqmaybe_swi BRel} & iff $\bot$ is aware under \code{BRel} & iff $\top$ is aware \\
\bottomrule
\end{tabularx}
\end{table}

\subsection{Compositional theorems}\label{sec:composition}

The calculus comes with one non-interference theorem per constructor
(\cref{tab:composition}). Each derives the equivalences of the composite from those
of its parts, which is what gives the equivalences of a large model their shape.
Two carry substantial side conditions. \m{sta_NI} asks that the state update map
related states to related states,
\[
\mathit{fv\_NI}:\quad \rel\;l\;i\;i' \wedge \rel\;l\;v\;v' \;\Rightarrow\;
  \rel\;l\;(f\;i\;v)\;(f\;i'\;v'),
\]
and that an unobservable input not visibly move the state,
$\mathit{f\_EP}:\; \dis\;l\;i \Rightarrow \rel\;l\;(f\;i\;v)\;v$. \m{swi_NI} asks, at
each level, that the observer either can trust the gate bit (it is \emph{aware}) or
learns nothing from the gated process because all its outputs fall in one class (it
is \emph{oblivious}).

\begin{table*}[t]
\centering
\caption{The composition theorems, slightly simplified. $I$ and $O$ are the input
  and output equivalences, $V$ the state equivalence and $B$ the gate equivalence.}
\label{tab:composition}
\footnotesize
\begin{tabularx}{\textwidth}{@{}lLL@{}}
\toprule
theorem & from & to \\
\midrule
\m{out_NI} & --- & $\NI\;I\;O\;(\mathsf{out}\;o)$ \\
\m{map_NI} & $\NI\;I'\;O\;p$, side conditions on $f$ and $g$
  & $\NI\;I\;O'\;(\mathsf{map}\;f\;g\;p)$ \\
\m{sta_NI} & $\NI\;(\mathsf{eqpair\_R}\;V\;I)\;O\;p$, $\mathit{fv\_NI}$,
  $\mathit{f\_EP}$
  & $\NI\;I\;(\mathsf{eqpair}\;V\;O)\;(\mathsf{sta}\;f\;g\;v\;p)$ \\
\m{swi_NI} & $\NI\;I\;(\mathsf{eqpair\_LR}\;B\;O)\;p$, aware or oblivious at
  each level
  & $\NI\;(\mathsf{eqpair\_LR}\;B\;I)$\newline
    \hspace*{1em}$(\mathsf{eqmaybe\_swi}\;O\;B)\;(\mathsf{swi}\;b\;p)$ \\
\m{par_NI} & $\NI\;I\;O_1\;p_1$ and $\NI\;I\;O_2\;p_2$
  & $\NI\;I\;(\mathsf{eqpair}\;O_1\;O_2)\;(\mathsf{par}\;p_1\;p_2)$ \\
\m{loop_NI} & $\NI\;I\;I\;p$ & $\NI\;I\;I\;(\mathsf{loop}\;p)$ \\
\m{maybe_NI} & $\NI\;I\;O\;p$
  & $\NI\;(\mathsf{eqmaybe\_hidden}\;I)\;O\;(\mathsf{maybe}\;p)$ \\
\bottomrule
\end{tabularx}
\end{table*}

\section{An Operating System as One Process}\label{sec:model}

\Cref{sec:background} gave us a calculus of processes, equivalences that say what
an observer at each level sees, and one composition theorem per constructor. We now
use the calculus to write an operating system with interrupt handling as a single
process, following the example of \cref{fig:leak}: the secret is whether a disk
interrupt occurred. We first introduce the processes the system runs
(\cref{sec:pool}) and how they run together (\cref{sec:running}). That explains what
the system receives (\cref{sec:interface}) and what it emits
(\cref{sec:outputs}). We then give the rule that decides which process runs
(\cref{sec:state}). What is left open is the \emph{interrupt policy}
(\cref{sec:parameters}); \cref{sec:designs} gives two.

\subsection{The processes}\label{sec:pool}

The system runs six processes: three interrupt handlers, a scheduler, and two user
processes.

\paragraph{Handlers} The timer, disk and NOP handlers are copies of one process. A
handler run takes exactly \code{runtime} steps, and on its last step the handler
emits \code{Notify} to announce completion.

\paragraph{User processes} The public user process issues requests (\code{Get})
whose results the attacker sees. The private user process takes the disk handler's
output as input, and emits \code{Sys} when it has been notified of a disk event and
\code{Idle} otherwise.

\paragraph{Scheduler} When the scheduler runs, it proposes which user process runs
next.

\subsection{How the processes run together}\label{sec:running}

The six processes run side by side in a \emph{pool} (\cref{fig:system}), and only
one of them takes a step at a time. Each sits behind a switch, controlled by the
current process identifier (pid): a pid names one slot, whose process takes one
step, while every other slot emits \code{None} without moving. The pid is held in a
global state, which a state cell and a feedback loop thread around the pool.

\begin{figure}[t]
\centering
\begin{tikzpicture}[>=stealth, font=\scriptsize,
  slot/.style={draw, minimum width=2.3cm, minimum height=0.42cm, inner sep=0pt},
  outv/.style={minimum width=1.0cm, minimum height=0.42cm, inner sep=0pt,
    font=\ttfamily\scriptsize}]
% six slots and their outputs, on a common grid
\foreach \n/\l/\o [count=\k from 0] in
  {t/timer handler/None, d/disk handler/None, z/NOP handler/None,
   s/scheduler/None, v/private user/None, u/public user/Some o}
  {\node[slot] (\n) at (3.0, 1.05-0.42*\k) {\l};
   \node[outv] (o\n) at (5.15, 1.05-0.42*\k) {\o};
   \draw[->] (\n.east) -- (o\n.west);}
\node[draw, dashed, rounded corners=2pt, fit=(t)(u), inner sep=3pt] (pool) {};
\node[draw, rounded corners=2pt, fit=(ot)(ou), inner sep=2pt] (tuple) {};
\node[anchor=south] at (pool.north) {pool: six slots in \texttt{par}};
\node[anchor=south] at (tuple.north |- pool.north) {six-tuple};
% global state, vertically centred on the pool
\node[draw, rounded corners=2pt, minimum width=1.3cm, minimum height=1.1cm,
  align=center] (sta) at (0.35, 0 |- pool.center) {global\\state};
\draw[->] (-1.25, 0 |- sta.center) -- node[above] {interrupt} (sta.west);
\draw[->] (sta.east) -- node[above] {pid} (pool.west |- sta.east);
\draw[->] (tuple.east) -- ++(0.5,0) node[right] {out};
\draw[->] ($(tuple.east)+(0.25,0)$) |- ($(pool.south)+(0,-0.35)$) -|
  node[pos=0.25, below] {fed back (\texttt{loop})} (sta.south);
\end{tikzpicture}
\caption{The model. An interrupt only updates the global state. The state's pid
  selects one slot, which takes a step and emits \texttt{Some o}; every other slot
  emits \texttt{None}. Because the slots are composed with \texttt{par}, their
  outputs form a six-tuple, which is emitted and fed back to update the state.}
\label{fig:system}
\end{figure}

\subsection{Inputs}\label{sec:interface}

The system receives interrupts. Its input type \code{T_in} has two values, the timer
interrupt and the disk interrupt. An interrupt goes to the global state, where it
marks itself pending in the interrupt controller. The controller has one more
entry, for the NOP handler, which no external interrupt raises; the sliced policy of
\cref{sec:sliced} sets it pending itself, to use the NOP handler as filler.

The disk interrupt is the secret. The input equivalence \m{in_equiv} makes it, and
only it, unobservable to the attacker; otherwise interrupts are related only when
equal:
\begin{align*}
\mathit{ir} \doteq_\ell \unobs &\iff \mathit{ir} = \mathsf{Disk} \wedge \ell = \Lo, \\
\mathit{ir} \doteq_\ell \mathit{ir}' &\iff \mathit{ir} = \mathit{ir}' \vee
  (\mathit{ir} \doteq_\ell \unobs \wedge \mathit{ir}' \doteq_\ell \unobs).
\end{align*}

\subsection{Outputs}\label{sec:outputs}

By the rule for \code{par} (\cref{fig:semantics}), the pool emits the
outputs of all its slots together. So the system emits a six-tuple of options, the output
type \code{T_out}, with exactly one \code{Some} on each step. The output therefore
records not only what was computed but also \emph{who ran}: the schedule is part of
what the system emits. \Cref{fig:slots} shows what we intend to hide from the
attacker in each position. The exact output equivalence, and why the proof forces
its shape, come in \cref{sec:equivs}.

\begin{figure}[t]
\centering
\begin{tikzpicture}[font=\scriptsize,
  proc/.style={draw, rounded corners=1.5pt, minimum width=4.1cm,
    minimum height=0.46cm, inner sep=1pt},
  % indented on both sides like the nested boxes of the key
  hid/.style={proc, fill=white},
  ran/.style={proc, fill=black!10, minimum width=3.7cm},
  val/.style={proc, fill=black!28, minimum width=3.3cm}]
% the processes, in the order of Fig.~3, shaded by what the attacker learns
\foreach \s/\k/\n/\t in
  {val/0/timer handler/{\Proc{\emptyset}{\mathit{Hout}}},
   hid/1/disk handler/{\Proc{\emptyset}{\mathit{Hout}}},
   hid/2/NOP handler/{\Proc{\emptyset}{\mathit{Hout}}},
   val/3/scheduler/{\Proc{\emptyset}{\mathbb{N}}},
   ran/4/private user/{\Proc{\mathit{Hout}}{S}},
   val/5/public user/{\Proc{\emptyset}{G}}}
  {\node[\s] (p\k) at (0, -0.52*\k) {};
   \node[anchor=west] at (p\k.west) {\n};
   \node[anchor=east] at (p\k.east) {$\t$};}
% key: the shades as nested sets
\begin{scope}[shift={(2.55,-0.52)}]
\draw[rounded corners=3pt, fill=white] (0,0) rectangle (2.75,-1.6);
\draw[rounded corners=3pt, fill=black!10] (0.2,-0.42) rectangle (2.55,-1.48);
\draw[rounded corners=3pt, fill=black!28] (0.4,-0.86) rectangle (2.35,-1.36);
\node[anchor=north west] at (0,0) {attacker sees nothing};
\node[anchor=north west] at (0.2,-0.42) {sees that it ran};
\node[anchor=north west] at (0.4,-0.86) {\dots\ and its value};
\end{scope}
\end{tikzpicture}
\caption{The six processes, with their types, shaded by what the attacker learns
  of each (key on the right). $\mathit{Hout} = \{\mathtt{Notify},
  \mathtt{Nothing}\}$ is the handler output, which the private user reads from the
  disk handler; $S = \{\mathtt{Sys}, \mathtt{Idle}\}$, $G = \{\mathtt{Get}\}$, and
  the scheduler emits a pid.}
\label{fig:slots}
\end{figure}
\subsection{When does a handler run?}\label{sec:state}

Besides the current pid, the global state holds an interrupt controller with a
\emph{pending} bit and a \emph{mask} bit for each handler: timer, disk and NOP. A
handler can be started when its entry is pending and not masked; if several can,
the timer goes first, then the disk, then NOP. Handlers cannot be interrupted: starting one sets all masks.

The state changes once per event. On an interrupt, \m{set_pending} sets that
interrupt's pending bit. On a pool output, \m{check_scheduler} adopts the pid the scheduler
proposed, if the scheduler ran; then \code{service_policy}, part of the interrupt
policy (\cref{sec:parameters}), reacts to the output; and finally \m{initiate_next} decides
who runs next by a priority rule:
\begin{enumerate}
\item if a handler is running, it continues;
\item otherwise, if a handler can be started, it starts;
\item otherwise, if a handler has just finished, control returns to the user
  process it interrupted, or to the scheduler if a reschedule was requested;
\item otherwise, the current process continues.
\end{enumerate}
Just before the choice, \m{initiate_next} applies \m{restore_invariant}, which
\cref{sec:parameters} explains.
Appendix~\ref{app:model} gives the code.

This is enough to see the leak of \cref{fig:leak}. With the disk unmasked, a disk
interrupt sets its pending bit. At the next output, rule 2 applies and the disk
handler starts at once, and for the \code{runtime} steps of the handler the public
user's slot emits \code{None}.

\subsection{The interrupt policy}\label{sec:parameters}

Together, these parts form the generic model
\begin{lstlisting}
model runtime init_policy service_policy
      restore_invariant p_pub p_priv p_sched
  : Proc T_in T_out
\end{lstlisting}
(\m{model}, written out in Appendix~\ref{app:model}). Its first argument,
\code{runtime}, is the handler length, and its last three are the scheduler and the
two user processes. The pool, the layout of the state, \m{set_pending},
\m{check_scheduler} and \m{initiate_next} are fixed. The remaining three arguments
are where models differ:
\begin{itemize}
\item \code{init_policy}, the initial state. Its masks decide which interrupts can be
  serviced before anything has happened.
\item \code{service_policy}, which runs on each pool output, before
  \m{initiate_next}. It is the only place that sees a handler announce
  completion, and so the only place masks can be reopened: it decides when the next
  handler becomes eligible.
\item \m{restore_invariant}, which \m{initiate_next} applies just before choosing.
  It writes into the state an invariant that the other two already maintain, for
  the benefit of the proof (\cref{sec:restore}).
\end{itemize}
The first two decide \emph{when a handler may run}, and together we call them the
\emph{interrupt policy}. The third is not part of it: it never changes a state the
model reaches (\cref{thm:noop}), and it is the identity where the proof needs no
help. \Cref{tab:designs} gives all three for both policies.

\section{Two Designs}\label{sec:designs}

\Cref{tab:designs} sets the two designs side by side. Both are \m{model} at a
different triple, and that difference is the whole security story.

\begin{table*}[t]
\centering
\caption{The two designs.}
\label{tab:designs}
\footnotesize
\begin{tabularx}{\textwidth}{@{}lLL@{}}
\toprule
 & \m{model_immediate} & \m{model_sliced} \\
\midrule
\code{init} & all masks clear; no slice counter & every mask set except the timer's;
  counter at \code{Some 0} \\
\code{handler_preroutine} & on a handler's \code{Notify}, unmask everything; on the timer's,
  also request a reschedule & on the timer's \code{Notify}, open a slice and request a
  reschedule; at each boundary inside it, reopen the masks and queue a NOP run; tick
  the slice and close it \\
\m{restore_invariant} & \code{id} & re-imposes the slice invariant; the identity on every
  reachable state \\
a handler runs & as soon as its interrupt is serviced & only inside a slice opened
  by the public timer \\
non-interfering? & no (\cref{thm:leak}) & yes (\cref{thm:main}) \\
\bottomrule
\end{tabularx}
\end{table*}

\subsection{Ordinary interrupt handling, and the leak}\label{sec:leak}

\m{model_immediate} starts with every mask clear, so an interrupt is serviced the
moment it arrives. Its preroutine reacts to a handler announcing completion by
unmasking everything, and if the handler was the timer's, by requesting a
reschedule. Nothing constrains \emph{when} a handler runs.

The counterexample is stated for a concrete system: a public process that
repeatedly emits the request \code{GetRequest}, a private process that issues a system
call when the disk handler has notified it and \code{NOP} otherwise, a round-robin
scheduler, and handlers of length two.

\begin{theorem}[\m{model_immediate_not_NI}]\label{thm:leak}
\begin{multline*}
\neg\,\NI\;\texttt{in\_equiv}\;\texttt{out\_equivC}\\
\texttt{model\_immediate\_concrete}.
\end{multline*}
\end{theorem}

\begin{proof}[Proof sketch]
Start from the trace of two public requests, which the model admits. Insert a disk
interrupt at the front; it is unobservable at $\bot$, so non-interference demands
that the result still be a trace. It is not. The interrupt sets the disk pending bit,
the first request still goes through, and the following step belongs to the disk
handler, so the public slot carries \code{None} where the second request was. Refuting
this requires showing that a trace is \emph{not} admitted, which needs reductions to
be invertible (\cref{sec:mechanisation}).
\end{proof}

\subsection{Sliced interrupt handling}\label{sec:sliced}

\m{model_sliced} keeps the disk and NOP handlers masked outside a slice. Its
preroutine is a composition of three stages, read right to left:
\begin{lstlisting}
sliced_preroutine o =
  check_ir_count (@$\circ$@) check_handler_completed
  (@$\circ$@) initiate_ir o
\end{lstlisting}
\begin{itemize}
\item \emph{Opening a slice.} When the timer handler emits \code{Notify}, \m{initiate_ir} loads
  the counter with $\mathit{runs}\cdot\mathit{runtime}$, so that exactly \code{runs}
  handler runs fit in the slice, and requests a reschedule for when the slice ends.
\item \emph{Boundaries.} Because the slice length is a multiple of \code{runtime}, the counter
  is a nonzero multiple of \code{runtime} exactly when a handler run has just finished.
  There \m{check_handler_completed} reopens the disk and NOP masks, keeps the timer
  masked, and makes the NOP handler pending. That last step is the \emph{filler}: a
  handler is always ready at a boundary, so a run takes place whether or not a disk
  interrupt is waiting, and the disk handler, if pending, takes priority.
\item \emph{Closing.} \m{check_ir_count} ticks the counter, and when it reaches zero masks the
  disk and NOP handlers, unmasks the timer and disables the counter.
\end{itemize}

A secret handler run therefore takes the place of a NOP run instead of displacing a
user process: a slice containing a disk run has the same length and position as one
containing only filler, and the public schedule around it is unchanged.
\Cref{fig:traces} shows two runs of the concrete system. In the run with a disk
interrupt, the interrupt is merely recorded until the next slice opens, and the only
structural difference between the runs is which handler occupies the first run of
the slice.

\begin{figure*}[t]
\centering
\footnotesize
\begin{tabular}{@{}lll@{}}
\toprule
no disk interrupt & with disk interrupt & \\
\midrule
\code{Get} & \code{Get} & \\
 & \textit{disk interrupt} & recorded as pending; nothing else happens \\
 & \code{Get}, \code{Get} & \\
\textit{timer interrupt} & \textit{timer interrupt} & \\
\code{Get} & \code{Get} & \\
\code{tmr Nth}, \code{tmr Nfy} & \code{tmr Nth}, \code{tmr Nfy} & timer handler completes: slice opens \\
\code{dfl Nth}, \code{dfl Nfy} & \code{dsk Nth}, \code{dsk Nfy} & the disk handler replaces a NOP run \\
\code{dfl Nth}, \code{dfl Nfy} & \code{dfl Nth}, \code{dfl Nfy} & \\
\code{sch hi} & \code{sch hi} & slice over: the scheduler runs \\
\code{NOP} & \code{Sys} & the private process reacts; both are secret \\
\bottomrule
\end{tabular}
\caption{Two runs of the concrete sliced system (abridged). Handler slots are
  secret including whether they ran, so to the attacker both columns look alike.}
\label{fig:traces}
\end{figure*}

\m{model_sliced} still exposes every slot of the pool. The model
\m{model_sliced_userview} applies a projection \m{parse_output} that keeps only the two
user processes' outputs; it is \m{model_sliced} behind a \code{map}, and its
non-interference follows by weakening the output equivalence.

\section{The Security Proof}\label{sec:proof}

\begin{theorem}[\m{model_sliced_NI}]\label{thm:main}
If the scheduler and the public user process are non-interfering at
\m{public_equiv} on input and output, and the private user process is
non-interfering at \m{private_equiv} on input and output, then
\begin{multline*}
\mathsf{model\_sliced}\;\mathit{runtime}\;\mathit{runs}\;p_{pub}\;p_{priv}\;p_{sched}\\
  \in \NI(\mathsf{in\_equiv},\, \mathsf{out\_equiv}).
\end{multline*}
\end{theorem}

\begin{corollary}[\m{model_sliced_userview_NI}]
Under the same hypotheses, \m{model_sliced_userview} is non-interfering at
\m{out_equiv_userview}.
\end{corollary}

The corollary follows from \cref{thm:main} by \m{map_NI}, since \m{parse_output} maps
\m{out_equiv}-related outputs to \m{out_equiv_userview}-related ones. The rest of this
section is about \cref{thm:main}.

\subsection{The equivalences}\label{sec:equivs}

\paragraph{Input} \m{in_equiv} makes the disk interrupt, and only the disk interrupt,
unobservable at $\Lo$. The timer stays public, because the sliced policy schedules by it.

\paragraph{Output} A pool output is a six-tuple of options, one per slot, and
\m{out_equiv} is the corresponding nest of \m{eqpair}s (\cref{tab:out-equiv}). Its two
columns come apart only for the two secret handlers, and that is where scheduling
secrecy lives. The payload equivalence hides a slot's \emph{value}. Whether the slot
\emph{ran} is decided by the \code{eqmaybe} instance around it, which fixes who can see
\code{None}. Under \m{eqmaybe_all} the attacker tells \code{Some} from \code{None}, and so
sees that the private user process ran; under \m{eqmaybe_nonbot} a handler run and
a step in which nothing happened are the same observation.

\begin{table}[t]
\centering
\caption{\texttt{out\_equiv}, one component per slot: a wrapper from the
  \texttt{eqmaybe} family around a public or private payload equivalence. The last
  two columns say what the attacker at $\Lo$ sees.}
\label{tab:out-equiv}
\footnotesize
\begin{tabular}{@{}lllll@{}}
\toprule
slot & wrapper & payload & value & ran? \\
\midrule
public user & \m{eqmaybe_all} & public & visible & visible \\
private user & \m{eqmaybe_all} & private & hidden & visible \\
scheduler & \m{eqmaybe_all} & public & visible & visible \\
NOP & \m{eqmaybe_nonbot} & private & hidden & hidden \\
disk & \m{eqmaybe_nonbot} & private & hidden & hidden \\
timer & \m{eqmaybe_all} & public & visible & visible \\
\bottomrule
\end{tabular}
\end{table}

The NOP slot must be classified like the disk slot: a slice run is filled by one or
the other, so the two must look alike.

\subsection{Compositional structure}

The proof follows the term of \m{model_sliced} from the outside in, applying one
composition theorem per constructor (\cref{thm:composition}): \m{map_NI} for the outer
rewiring, \m{loop_NI} for the feedback, \m{sta_NI} for the global state, \m{maybe_NI} for
the idling pool, \m{par_NI} to split the pool into slots, \m{swi_NI} for each slot's
gate, and at the leaves \m{out_NI} for the handler processes and the hypotheses of
\cref{thm:main} for the scheduler and user processes. Since each theorem derives the
composite's equivalences from those of its parts, the shape of \m{out_equiv} is forced:
\m{par_NI} imposes its nest of \m{eqpair}s and \m{swi_NI} its \code{eqmaybe} instances.

Two of these steps carry substantial side conditions: the gate obligations of
\m{swi_NI} (\cref{sec:gating}) and the state obligations of \m{sta_NI}
(\cref{sec:state-equiv,sec:fvni}).

\subsection{Gating a slot}\label{sec:gating}

Each slot is a \code{swi} whose gate bit is computed from the pid by a \code{map}. Two
obligations remain per slot, both parameterised by an equivalence $\eqv{B}$ on the gate
bit that the proof may choose. \m{map_NI} requires the gate map to be unobservable-preserving,
sending unobservable pids to unobservable bits (\m{f_PU}), and \m{swi_NI} requires, at every level, that the
observer be aware of the gate or oblivious to the slot:
\[
\ell \in \Aw(\mathsf{true}, \eqv{B}) \;\cup\;
  \Ob(\mathsf{swi}\;b\;p,\; \mathsf{eqmaybe\_aware}(\eqv{O},\eqv{B})).
\]
Awareness means a bit related to \code{true} is \code{true} and observable, so the observer can
trust the gate. Obliviousness means there is a reference output $o_0$ such that every
output of the slot is related to it: whatever the input does, the observer sees the
same class. This is weaker than the original condition (\cref{sec:departures}), and the
argument below needs the weaker form.

At $\Lo$ the two pids of the secret handlers are unobservable. For a slot neither of
them selects, both map to \code{false}, so a $\eqv{B}$ in which \code{false} alone is unobservable
satisfies \m{f_PU}; we call it \m{false_equiv}. Under \m{false_equiv} the observer is aware
at every level, and the obligation is discharged without a case split. This covers
the public and private user slots, the scheduler and the timer handler.

For the disk and NOP slots, one secret pid maps to \code{true} and the other to \code{false},
so \m{f_PU} demands that both bits be unobservable, which leaves only \m{private_equiv}.
But then \code{true} is unobservable at $\Lo$ and awareness fails there. At $\Hi$
awareness still holds; at $\Lo$ the proof uses obliviousness, with \code{None} as the
reference output. A step where the slot is gated off emits \code{None}, and a step where
it runs emits \code{Some o}, so the obligation is that \code{None} and \code{Some o} are related at
$\Lo$. Under \m{eqmaybe_aware} that holds exactly when the gate is hidden and the
payload unobservable, which the \m{private_equiv} payload of \cref{tab:out-equiv}
provides. The line therefore falls between the two secret handlers and everything
else; the private user process sits on the public side, its secrecy carried by its
payload and never by its gate.

\subsection{The state equivalence}\label{sec:state-equiv}

\m{sta_NI} imposes two conditions on the state update, and they pull in opposite
directions.
\begin{itemize}
\item \m{f_EP} requires that an unobservable input not visibly move the state. A disk
  interrupt is unobservable at $\Lo$ and sets the disk pending bit, so that bit
  must be private. (The event equivalence \code{eqsum_L in_equiv out_equiv} makes only
  interrupts unobservable, so \m{f_EP} constrains \m{set_pending} alone.)
\item \m{fv_NI} requires that related states be mapped to related states.
  \m{initiate_next} branches on the mask bits, so two related states that disagreed
  there would take different branches; the mask bits must be public.
\end{itemize}
Both hold only if the interrupt controller is classified per bit: the disk and NOP
pending bits are private and every mask bit is public (\cref{tab:state-equiv}).
The same constraint explains why \m{set_pending} sets one bit and stops: any decision
taken on an input step, such as starting a handler or moving the slice counter,
would disturb a public field. Everything an interrupt causes is deferred to a later
output step, where the input is no longer the thing being varied.

\begin{table}[t]
\centering
\caption{\texttt{state\_type\_equiv}, field by field.}
\label{tab:state-equiv}
\footnotesize
\begin{tabularx}{\columnwidth}{@{}>{\raggedright\arraybackslash}p{0.25\columnwidth}L>{\raggedright\arraybackslash}p{0.24\columnwidth}@{}}
\toprule
field & equivalence & at $\Lo$ \\
\midrule
current pid & \code{eqsum private_equiv public_equiv} & which secret handler is hidden \\
previous pid, \m{re_sched}, \m{ir_count} & \m{public_equiv} & visible \\
NOP and disk bits & \m{private_pending_equiv} & pending hidden, mask visible \\
timer bits & \m{public_bits_equiv} & visible \\
\bottomrule
\end{tabularx}
\end{table}

The pid type is $\mathsf{Sum}\;\mathsf{Bool}\;\mathsf{Nat}$. Its \code{inl} side
holds exactly the two secret handlers, the disk handler and the NOP handler, and its
\code{inr} side everything else, the timer handler included. Drawing the line there
lets the proof classify a whole pid in one place.

The current pid is where the policy does its work. Since \m{eqsum} never relates an
\code{inl} to an \code{inr}, the fact \emph{that} a handler is running stays visible, as does
every public pid. What \m{private_equiv} hides is the bit inside \code{inl}: at $\Lo$ the
disk handler and the NOP handler look alike. Across two $\Lo$-related states, that
bit and the two private pending bits are the only things allowed to differ.

\subsection{Closure under the transition}\label{sec:fvni}

The core obligation is that \m{state_type_equiv} is closed under the whole transition.
The transition is a sum, \m{set_pending} on interrupts and a pipeline of stages on
outputs, and the event equivalence never relates an interrupt to an output. So the
goal splits along the sum, and a composition lemma (\m{fv_NI_comp}) breaks the
pipeline into one goal per stage:
\begin{itemize}
\item \m{set_pending}: related interrupts set related pending bits;
\item \m{check_scheduler}: the scheduler's pid is public, so the update agrees;
\item \code{service_policy}: the slice logic reads only the public counter and
  masks, and writes the same values into both states;
\item \m{initiate_next}: related states must make the same choice of what runs
  next.
\end{itemize}
Each goal is self-contained, which is the point of splitting. The cost is that a
stage must be proved for \emph{all} pairs of related states, since it cannot assume
its input came from the stage before. That loses the invariant of the time slice,
and the last stage cannot be proved without it.

\subsection{Restoring the invariant}\label{sec:restore}

Two facts hold of every state \m{model_sliced} reaches. While a slice is live and no
handler is running, the NOP handler is pending, the disk and NOP masks are clear and
the timer is masked; once the slice is over, the disk and NOP handlers are masked.
And the disk and NOP masks always agree. The sliced policy maintains both
(\cref{sec:sliced}).

The function \m{restore_invariant} writes these facts into the state, just before
\m{initiate_next} chooses. It reads only the public counter, so it sets the same bits in
two related states, and afterwards the deciding mask bits agree and related states
take the same branch. This is what discharges the \m{initiate_next} goal.

We prove that it never changes a state the model reaches.

\begin{theorem}[\m{restore_invariant_noop}]\label{thm:noop}
Let $\mathit{runtime} \neq 1$. For every sequence of events $es$ from the initial
state of \m{model_sliced}, and every output $o$, let $w$ be the state that
\m{initiate_next} receives on $o$ after $es$. If not all masks are set in $w$, then
$\mathsf{restore\_invariant}\;w = w$.
\end{theorem}

The events $es$ are arbitrary, so the theorem holds whatever the pool emits. The proof states the invariant
between steps, shows that every event preserves it, and shows that at the choice
point a slightly weaker invariant holds, on which \m{restore_invariant} is the
identity. One small lemma per stage of the pipeline does the bookkeeping.

The hypothesis $\mathit{runtime} \neq 1$ exposes a quirk of the model. The boundary
test reads the counter before the stage that ticks it, while \m{restore_invariant}
reads it after. With one-step handlers the last boundary inside a slice ticks the
counter to zero in the same step, so the slice already reads as over, and
\m{restore_invariant} masks the handler that boundary had just enabled. Non-interference
is unaffected (\cref{thm:main} has no such hypothesis), but at $\mathit{runtime} = 1$
every slice loses its last handler run.

An earlier version of the model relied on \m{restore_invariant} for behaviour: it was
the only thing that queued the NOP filler and the only thing that returned control
to the scheduler after a slice. Proving \cref{thm:noop} required moving both into
\code{service_policy}, where they now appear as part of the policy
(\cref{sec:sliced}). The proof of \cref{thm:main} needed only two small lemmas to
follow the change.

\section{How General the Result Is}\label{sec:generality}

\Cref{thm:main} is not about one system. Discharging the section variables of the
mechanisation, it reads in full:
\begin{lstlisting}
forall (Opub Opriv : Ty) (runtime runs : nat)
       (p_pub : Proc Empty Opub)
       (p_priv : Proc THandlerOutput Opriv)
       (p_sched : Proc Empty Nat),
  NI (public_equiv Empty) (public_equiv Opub)
     p_pub ->
  NI (private_equiv THandlerOutput)
     (private_equiv Opriv) p_priv ->
  NI (public_equiv Empty) (public_equiv Nat)
     p_sched ->
  NI in_equiv (out_equiv Opub Opriv)
     (model_sliced runtime runs
                   p_pub p_priv p_sched)
\end{lstlisting}
Everything the pool is built around stays fixed; everything it carries is left
open (\cref{tab:open}). The side conditions are the expected ones: the mechanism
cannot make a leaking process secure, so each slot is assumed non-interfering at its
own classification, and the theorem says that composing them through this mechanism
leaks nothing further. The handler length and slice size carry no side condition,
because the slice is \emph{defined} as $\mathit{runs}\cdot\mathit{runtime}$ and so
ends on a handler boundary by construction.

\begin{table}[t]
\centering
\caption{What \cref{thm:main} leaves open.}
\label{tab:open}
\footnotesize
\begin{tabularx}{\columnwidth}{@{}lLl@{}}
\toprule
parameter & ranges over & side condition \\
\midrule
\code{p_pub} & any public user process & $\NI$ at \m{public_equiv} \\
\code{p_priv} & any private user process & $\NI$ at \m{private_equiv} \\
\code{p_sched} & any scheduler & $\NI$ at \m{public_equiv} \\
\code{Opub}, \code{Opriv} & the user processes' output alphabets & none \\
\code{runtime} & handler length & none \\
\code{runs} & handler runs per slice & none \\
\bottomrule
\end{tabularx}
\end{table}

Two features of the design make this possible. The state update never reads the
\emph{value} a user slot emits, only whether it emitted (the scheduler-output test
inspects the user slots for \code{None} and nothing more), so no user behaviour reaches
the schedule. And the state obligation \m{fv_NI} never mentions the slot processes, so
the one hard obligation is unaffected by what fills them.

What stays fixed is the mechanism under study, the three handlers, and the shape of
the pool around it: six slots in a fixed order. The order is hard-wired in the
projections that read the pool's output, and varying the number of user slots would
change the state transition and with it the proof of \m{fv_NI}.

\paragraph{Levels} We have read levels as two. The mechanisation works over an
arbitrary bounded lattice, where a characterised equivalence must also be monotone
downwards (what a level relates stays related below it, and likewise for
unobservability), and $\NI$ quantifies over every level. Nothing in the proof depends
on the lattice having two points, because every equivalence in the models
distinguishes $\bot$ from all other levels and nothing more. The level case splits
in the proof are $l = \bot$ against $l \neq \bot$, and the second case covers every
other level of any lattice at once.

\section{Mechanisation}\label{sec:mechanisation}

The development is about 6{,}000 lines of Rocq 9.0.1, available at
\url{\repobase}, using Mathematical
Components~\cite{mathcomp} and, for the coinductive definitions of the original
paper, paco~\cite{paco} (\cref{tab:files}).

\begin{table}[t]
\centering
\caption{The development.}
\label{tab:files}
\footnotesize
\begin{tabularx}{\columnwidth}{@{}lrL@{}}
\toprule
file & lines & contents \\
\midrule
\mechfile{theories/definitions.v}{definitions.v} & 908 & the calculus, traces, $\NI$, the equivalences \\
\mechfile{theories/theorems.v}{theorems.v} & 1337 & the composition theorems \\
\mechfile{theories/models.v}{models.v} & 1079 & the generic model, both policies, \cref{thm:noop}, example traces \\
\mechfile{theories/noninterference.v}{noninterference.v} & 1749 & the model's equivalences, \cref{thm:leak,thm:main} \\
\mechfile{theories/adequacy.v}{adequacy.v} & 904 & the original definitions, the equivalences with ours, \cref{thm:leak,thm:main} in the original's terms \\
\bottomrule
\end{tabularx}
\end{table}

\paragraph{Types} Input and output types range over an inductive universe \m{Ty} of
codes, interpreted into \code{Set}, rather than over \code{Set} directly. This makes
reductions invertible, which the negative result needs: refuting non-interference
means showing that a trace is \emph{not} admitted.

\paragraph{Axioms} The development assumes nothing beyond the two parameters
that fix the security lattice: no classical logic, and no uniqueness of identity
proofs. The latter is what dependent elimination on the type-indexed reduction
relations would need; we invert reductions through determinism of reduction
instead, one lemma per rule (\m{inv_mapO} and its siblings).

\subsection{Non-interference over finite traces}\label{sec:departures}

Our definitions follow Rafnsson et al., with one essential departure. The original
defines non-interference as a coinductive simulation over infinite streams of
labels, whose clauses apply at the head of a stream. We quantify over finite traces
instead, with closure clauses that apply at an arbitrary position (Definition 1).
This makes every proof in the development an induction on a trace.

The original definition is written out in \mechfile{theories/adequacy.v}{adequacy.v},
and \m{NI_paper} proves the two equivalent. The proof passes through infinite
streams (\m{NI_SNI}) and rests on two properties: finite and infinite behaviour
determine each other, which needs reduction to be deterministic (to assemble
finite runs into one infinite run) and never to block (to continue a finite run);
and insertion commutes with taking prefixes. Obliviousness, the other definition
that quantifies over behaviour, is likewise stated over traces and equivalent to
the original (\m{ObliviousDis_iff}). As a consequence, \cref{thm:leak,thm:main}
also hold with non-interference read as in the original, hypotheses included
(\m{model_immediate_not_NI_paper}, \m{model_sliced_NI_paper}).

One further difference is only notational. The original adjoins the distinguished
value $\unobs$ to every type; we keep the values as they are and make
unobservability a predicate, tied to the equivalence by the characterisation law.
\m{NI_lEquiv} shows that the two give the same notion of non-interference, and
this paper uses the original notation throughout.

The composition theorems themselves are proved directly for our definitions.

\section{Related Work}\label{sec:related}

% TODO: this section is a starting point; extend before submission.

\paragraph{Non-interference} Non-interference goes back to Goguen and
Meseguer~\cite{goguen1982}; Sabelfeld and Myers survey language-based
approaches~\cite{sabelfeld2003}. Timing channels in programs have been closed by
padding transformations that equalise the running time of secret-dependent
branches~\cite{agat2000}; the sliced policy applies the same idea to handler
runs, with the NOP handler as padding.

\paragraph{Composable timing-sensitive non-interference} Our proofs rest on the
calculus of Rafnsson et al.~\cite{rafnsson2017}, whose definitions are
timing-sensitive by construction and whose theorems compose. We mechanise the
calculus and its composition theorems, and apply them to a model of interrupt
handling.

\paragraph{Verified operating systems} Machine-checked information-flow proofs
exist for real kernels, notably seL4~\cite{murray2013}, whose proof explicitly
leaves timing channels out of scope. Time protection~\cite{ge2019} argues that an
operating system must partition time as well as space, and includes the
partitioning of interrupts among its mechanisms\todo{check this characterisation
against~\cite{ge2019}}; our sliced policy is an instance of
that idea in a small model. What we add is a timing-sensitive, compositional,
mechanised proof that the mechanism is secure, for any non-interfering scheduler and
user processes.

\section{Limitations and Conclusion}\label{sec:conclusion}

\paragraph{Limitations} The model makes several simplifications.
\begin{itemize}
\item \emph{One secret interrupt.} The NOP handler exists only as the disk handler's
  indistinguishable partner. Several secret interrupts would need a richer
  indistinguishability argument.
\item \emph{Equal handler lengths.} Every handler runs for the same number of steps,
  which is what lets every mask bit be public. Variable-length handlers would force
  the masks private, and the proof of \m{fv_NI} would have to reason across diverging
  branches. A real system would reach a fixed length by padding.
\item \emph{Cooperative processes, no nesting.} A selected process takes exactly one
  step; preemption within a step is not modelled. Handlers cannot be interrupted,
  since all masks are set while one runs.
\item \emph{Abstract time.} One output step is one unit of time, for every process
  and handler alike. Caches, memory contention and other microarchitectural channels
  are outside the model.
\item \emph{One-step handlers.} At $\mathit{runtime} = 1$ a slice loses its last
  handler run (\cref{sec:restore}); the result still holds, but the policy behaves
  less usefully there.
\end{itemize}

\paragraph{Conclusion} A secret interrupt can leak through the schedule alone, and
servicing interrupts only in publicly timed, padded slices closes the channel. We
proved both facts in Rocq, the second for any non-interfering scheduler and user
processes, by composing generic non-interference theorems over the structure of the
model. The proof did more than check the policy: it fixed the classification of the
interrupt controller bit by bit, separated hiding what a handler computed from hiding
that it ran, and exposed a proof device that was quietly carrying behaviour, which we
moved into the policy and then proved inert.

% TODO: future work, e.g. several secret interrupts, variable-length handlers with
% padding, or connecting the model to a real interrupt controller.


\bibliographystyle{IEEEtran}
\bibliography{references}

\end{document}
